% Semantic Similarity Search Service — Technical Report
\documentclass[11pt,a4paper]{article}

% ── Layout ────────────────────────────────────────────────────────
\usepackage[margin=2.5cm]{geometry}
\usepackage{parskip}
\usepackage{graphicx}
\usepackage{xcolor}
\usepackage{url}

% ── Typography ────────────────────────────────────────────────────
\usepackage{fontspec}
\usepackage{amsmath}
\setmainfont{DejaVu Sans}
\setmonofont{DejaVu Sans Mono}

% ── References ────────────────────────────────────────────────────
\usepackage{hyperref}
\hypersetup{hidelinks}

% ── Misc ──────────────────────────────────────────────────────────
\usepackage{listings}
\lstset{basicstyle=\ttfamily\small, breaklines=true, aboveskip=6pt, belowskip=6pt}

\title{Semantic Similarity Search for Unfair Clause Detection\\\large Technical Report}
\author{S. N. Thibaud\\\small\texttt{snthibaud@gmail.com}}
\date{}

\begin{document}
\maketitle

\begin{center}\rule{0.5\linewidth}{0.5pt}\end{center}

% ══════════════════════════════════════════════════════════════════
\section{Problem \& Approach}

Terms of Service documents routinely contain clauses that may be unfair
to consumers.  The CLAUDETTE project~\cite{lippi2019claudette} provides
an annotated corpus of 9,414 sentences from 50 ToS documents (from 12,011
raw sentences after discarding fragments shorter than five words), each
labelled along eight unfairness categories.  Only 1,032 sentences (11\%)
carry any unfairness label.

The goal is a production-style service that indexes the corpus,
supports semantic similarity search, classifies unfair clauses,
and is containerised with Docker.  Generative language models and
retrieval-augmented generation are out of scope.

\textbf{Methodology.}  Each experiment tests a specific hypothesis,
produces a quantitative result, and raises the next question.  The
progression is:

\begin{enumerate}
\item \textbf{Baseline}: embed with Legal-BERT, classify with Logistic
  Regression.  Is this viable?
\item \textbf{Domain}: does legal-domain pre-training matter?  Compare
  against a general-domain SBERT model.
\item \textbf{SBERT fine-tuning}: can contrastive fine-tuning of the
  encoder improve results without changing the architecture?
\item \textbf{Significance}: are the observed differences real, or within
  statistical noise?  Bootstrap the ranking metric.
\item \textbf{External validation}: how does the system compare to Lippi
  et al.~\cite{lippi2019claudette} under their exact evaluation protocol?
\end{enumerate}

All code was developed iteratively with an AI coding assistant (pi) and
manually verified: every reported number can be reproduced with the
commands in Section~\ref{sec:running}.

\textbf{Metrics.}  Three metric families are used throughout:

\begin{description}
\item[MCC] Matthews Correlation Coefficient, for $N$ samples and
  category $c$:
  \[\text{MCC}_c =
    \frac{\text{TP}_c \cdot \text{TN}_c -
          \text{FP}_c \cdot \text{FN}_c}
         {\sqrt{(\text{TP}_c+\text{FP}_c)(\text{TP}_c+\text{FN}_c)
                (\text{TN}_c+\text{FP}_c)(\text{TN}_c+\text{FN}_c)}}\]
  Accounts for all four confusion-matrix quadrants, unbiased by class
  imbalance ($\sim$11\% positive).  Used for model selection ($C$
  tuning) and threshold calibration.  Requires a chosen threshold.
  Reported as macro-average across the 8 categories.
\item[AP] Average Precision, the area under the precision-recall curve
  for a ranked list of $N$ predictions with true labels $y_i \in
  \{0,1\}$:
  \[\text{AP} = \sum_{k=1}^{N} \text{Prec}(k) \cdot
    \Delta\text{Rec}(k)\]
  where $\text{Prec}(k)$ and $\text{Rec}(k)$ are precision and
  recall at rank $k$.  Threshold-agnostic; measures ranking quality
  across all operating points.  Reported with bootstrap 95\%
  confidence intervals (2,000 paired iterations).
\item[Search metrics] Semantic search is evaluated via label overlap as
  a weak-supervision proxy: two sentences are \emph{relevant} iff they
  share $\ge 1$ positive unfairness label.  Let $Q$ be the set of
  1,032 query sentences that carry $\ge 1$ label (fair sentences
  carry no evaluation signal).  For a query $q \in Q$, let $R_q$ be
  its top-$k$ retrieved neighbours:
  \[\text{Prec}@k = \frac{1}{|Q|} \sum_{q \in Q}
    \frac{|\{r \in R_q : \text{relevant}(q, r)\}|}{k}\]
  \[\text{MRR} = \frac{1}{|Q|} \sum_{q \in Q}
    \frac{1}{\text{rank}_q},\quad
    \text{rank}_q = \text{position of first relevant result (0 if none)}\]
  \[\text{Gain} = \frac{\text{Prec}@k}{\text{random-Prec}@k}\]
  where random-Prec@$k$ uses $k$ uniformly sampled non-query
  sentences as the baseline.  This proxy is a lower bound ---
  sentences can be similar without sharing a label --- but it
  requires no additional annotation.
\end{description}

% ══════════════════════════════════════════════════════════════════
\section{Experiments}

\subsection{Baseline: Legal-BERT + Logistic Regression}

\textbf{Hypothesis}: frozen embeddings from a domain-adapted BERT model,
classified by a simple linear model, can capture enough legal semantics
for both search and unfair-clause detection.

\textbf{Setup}: encode all 9,414 sentences with
\texttt{nlpaueb/legal-bert-base-uncased}~\cite{chalkidis2020legalbert}
(110M parameters, 768-d, mean pooling over the last hidden layer), store
in a FAISS \texttt{IndexFlatIP}~\cite{johnson2019faiss} with L2-normalised
vectors, and train a multi-output Logistic Regression
(\texttt{class\_weight='balanced'}, $C$ tuned via 5-fold CV) on the
frozen embeddings.  Evaluate on a stratified 85/15 validation split.

\textbf{Result}: Validation MCC = 0.720, macro AP = 0.736.  The classifier
is viable: per-category MCC ranges from 0.61 (Limitation of Liability) to
0.83 (Unilateral Change---Usage).  Precision@10 for semantic search is
0.565, a 24.6$\times$ gain over random.

\textbf{Next question}: how much does the legal-domain pre-training
contribute?

\subsection{Domain pre-training: Legal-BERT vs.\ all-MiniLM}

\textbf{Hypothesis}: domain-specific pre-training provides a measurable
advantage over a general-purpose embedding model.

\textbf{Setup}: replace the encoder with
\texttt{all-MiniLM-L6-v2} (22M parameters, 384-d), keeping the
pipeline otherwise identical.

\textbf{Result}: all-MiniLM achieves MCC = 0.651 and AP = 0.657 --- a drop
of 7 MCC points and 8 AP points.  Search Precision@10 falls from 0.565 to
0.352.  The legal-domain advantage is substantial and consistent across
metrics.

\textbf{Next question}: the base Legal-BERT was pre-trained only with
masked language modelling.  Can SBERT-style contrastive fine-tuning
further improve the embedding quality?

\subsection{SBERT fine-tuning: Legal-BERT-Embedding}

\textbf{Hypothesis}: fine-tuning the encoder with a contrastive (Siamese)
objective produces cosine-optimised sentence embeddings that improve
both search and classification, without changing the inference
architecture.

\textbf{Setup}: add \texttt{maticzav/legal-bert-embedding} --- the same
Legal-BERT backbone fine-tuned as a sentence-transformers model --- to the
model registry.  Same pipeline, same 768 dimensions, same inference cost.

\textbf{Result}:

\begin{center}
\begin{tabular}{l c c c c c}
\hline
Model & Params & Dim & Build (CPU) & Val MCC & Val AP \\
\hline
Legal-BERT               & 110M & 768 & $\sim$21\,min & 0.720 & 0.736 \\
Legal-BERT-Embedding     & 110M & 768 & $\sim$28\,min & \textbf{0.764} & \textbf{0.742} \\
all-MiniLM-L6-v2         &  22M & 384 & $\sim$1.5\,min & 0.651 & 0.657 \\
\hline
\end{tabular}
\end{center}

Legal-BERT-Embedding improves MCC by 4.4 points and search Precision@10
from 0.565 to 0.617 (a 26.9$\times$ gain over random).  However, the AP
gain is only 0.6 points (0.736 $\rightarrow$ 0.742).

\textbf{Next question}: is the 0.6-point AP improvement statistically
significant, or just noise?

\subsection{Significance: bootstrap analysis}

\textbf{Hypothesis}: the SBERT fine-tuning genuinely improves ranking
quality (AP), not just threshold-dependent metrics like MCC.

\textbf{Setup}: paired bootstrap (2,000 iterations, 95\% CI) on the macro
AP difference between model pairs.

\textbf{Result}: the AP difference between Legal-BERT-Embedding and base
Legal-BERT is $\Delta = +0.006$, CI $[-0.038, +0.061]$ --- the interval
crosses zero, so the difference is \textbf{not statistically significant}.
The MCC gain is therefore driven by better probability calibration
enabling superior threshold selection, not by fundamentally better
ranking.  In contrast, the AP gap between either legal-domain model and
all-MiniLM ($\sim$0.08, CI $[+0.019, +0.158]$) \textbf{is} significant.

\textbf{Next question}: how does the system compare to the published
CLAUDETTE results~\cite{lippi2019claudette}?

\subsection{External validation: leave-one-document-out}

\textbf{Hypothesis}: under the same evaluation protocol used by Lippi et
al.\ (50-fold leave-one-document-out), our frozen-embedding approach is
competitive with their hand-crafted feature ensemble.

\textbf{Setup}: replicate the LOO protocol: for each of the 50 documents,
train on the other 49 (with fixed $C=50$; our earlier sweep shows MCC
varies by less than 0.01 for all $C > 10$), test on the held-out
document.  Train 8 independent Logistic Regressions on unfair-only
sentences for categorization, matching Lippi's architectural setup.

\textbf{Result}:

\begin{center}
\begin{tabular}{l c c}
\hline
 & Detection F1 & Categorization macro F1 \\
\hline
Lippi et al.\ (ensemble / 8 SVMs)       & 0.806 & $\sim$0.88 \\
Our Legal-BERT + multi-LR / 8-LR        & 0.742 / 0.753 & 0.889 \\
\hline
\end{tabular}
\end{center}

On detection, our single multi-output classifier (trained on all 9,414
sentences, fair and unfair) trails their 5-model ensemble (F1 0.75 vs.\
0.81).  This gap is partly due to model capacity (one linear model vs.\
an ensemble with complementary feature types) and partly due to the
joint detection+categorisation task being harder than detection alone.
On categorization, where the 8 independent classifiers are trained
\textbf{only on unfair sentences} (matching Lippi's setup: each classifier
learns to distinguish its category from other unfair categories, never
seeing fair sentences), our approach is competitive (0.889 vs.\ $\sim$0.88)
despite being dramatically simpler (linear classifiers on dense
embeddings vs.\ SVMs with hand-crafted features).  Per-category F1 ranges
from 0.77 (Content Removal) to 0.94 (Choice of Law, Jurisdiction).

% ══════════════════════════════════════════════════════════════════
\section{System Architecture}

Informed by the experiments above, the service uses Legal-BERT-Embedding
as the default encoder (best MCC and search quality at identical inference
cost to base Legal-BERT; all-MiniLM is available for faster builds).

\textbf{Components.}  The encoder produces a 768-d vector per sentence via
mean pooling.  FAISS \texttt{IndexFlatIP} stores all 9,414 vectors (exact
search, $\sim$1--2\,ms on CPU).  A multi-output Logistic Regression
classifier ($C=50$, \texttt{class\_weight='balanced'}) with per-category
MCC-optimal thresholds predicts the eight unfairness categories from the
same frozen embeddings.  (The production system is single-stage: one
classifier trained on all sentences handles both detection and
categorisation simultaneously.  For the Lippi comparison in \S2.5 we also
train 8 independent classifiers on unfair-only sentences --- same model
architecture, harder training data --- to match their two-stage protocol.)

\textbf{API.}  FastAPI with four endpoints: \texttt{/health} (GET),
\texttt{/search} (POST), \texttt{/classify} (POST), \texttt{/analyze}
(POST, combined).  Pydantic validation (1--2000 characters, \texttt{top\_k}
1--50).  Models load once at startup; no network calls at inference.

\textbf{Key design decisions.}

\begin{description}
\item[Shared frozen encoder.]  Using the same encoder for search and
  classification eliminates an extra model dependency.  The frozen
  weights make the classifier embedding-agnostic --- swapping encoders
  requires only rebuilding the FAISS index.
\item[Mean pooling over [CLS].]  Reimers \& Gurevych~\cite{reimers2019sbert}
  showed mean pooling of hidden states produces substantially better
  sentence embeddings than the [CLS] token for semantic similarity.
\item[Exact search.]  At 9,414 vectors, \texttt{IndexFlatIP} is both fast
  and precise.  Approximate indices are unnecessary at this scale but
  remain a one-line change for larger corpora.
\item[MCC as primary metric.]  MCC accounts for all four confusion-matrix
  quadrants and is unbiased by class imbalance ($\sim$11\% positive).
  AP (threshold-agnostic) with bootstrap CIs is used to assess whether
  improvements are genuine or threshold artifacts.  F1 is reported only
  for Lippi comparability.
\end{description}

\textbf{Testing.}  Property-based tests (Hypothesis~\cite{hypothesis})
verify API invariants ($\forall$ text, \texttt{/classify} returns exactly
8 categories; all scores $\in [0,1]$).  715 test cases pass.  Conventional
pytest covers the remaining cases.

% ══════════════════════════════════════════════════════════════════
\section{Limitations \& Trade-offs}

\begin{description}
\item[Frozen embeddings.]  Keeping the encoder frozen enables direct model
  comparison but leaves task-specific structure unexploited.  Contrastive
  fine-tuning with the eight labels could further improve both search and
  classification.
\item[Evaluation protocol.]  Our LOO implementation uses a fixed $C=50$
  rather than re-tuning per fold.  The MCC plateau above $C=10$ makes this
  harmless in practice, but a fully nested protocol would be more rigorous.
  The categorization LOO was run only for Legal-BERT; cross-model LOO
  comparison remains future work.
\item[Class imbalance.]  At $\sim$11\% positive and per-category rates as
  low as 2--3\%, rare categories (especially Arbitration) show high
  variance.  We mitigate with balanced class weights,
  stratified splits, and per-category MCC-optimal thresholds.
\item[Production gaps.]  The API is stateless and validates inputs but
  lacks authentication, rate limiting, and persistent logging.
\end{description}

% ══════════════════════════════════════════════════════════════════
\section{Running the Service}\label{sec:running}

\textbf{Docker (recommended).}  The service and UI run with a single
command:

\begin{lstlisting}[language=bash]
docker compose up --build
\end{lstlisting}

This builds Legal-BERT-Embedding inside the image ($\sim$28\,min on CPU,
$\sim$1.5\,min with \texttt{all-minilm}).  To use a different encoder:

\begin{lstlisting}[language=bash]
docker compose build --build-arg MODEL=all-minilm
docker compose up
\end{lstlisting}

The API listens on port 8000, the Streamlit UI on port 8501.

\textbf{Development (uv\footnote{\url{https://docs.astral.sh/uv}}).}  For local development or reproducing individual
experiments:

\begin{lstlisting}[language=bash]
# Install
uv sync --extra dev

# Build and train (default: legal-bert-embedding, ~28 min on CPU)
uv run python -m src.build_models
uv run python -m src.train_classifier

# Alternative encoders
uv run python -m src.build_models --model all-minilm          # ~1.5 min
uv run python -m src.build_models --model legal-bert-embedding  # ~28 min
uv run python -m src.train_classifier --model <model-key>

# Evaluation scripts (reproduce every number in this report)
uv run python -m src.evaluate_search --model <model-key>
uv run python -m src.bootstrap_ap --n-bootstrap 2000
uv run python -m src.evaluate_loo --model legal-bert --categorization

# Service + UI
uv run python -m src.service       # API on port 8000
uv run python -m src.ui            # Streamlit on port 8501
\end{lstlisting}

% ══════════════════════════════════════════════════════════════════

\bibliographystyle{plain}
\nocite{*}
\bibliography{references}

\end{document}
